\section{VC Dimension}

The homework and final materials repeatedly use shattering, growth functions,
and VC dimension to move from finite hypothesis classes to infinite classes.

\begin{definition}[Shattering]
A concept class $\C$ shatters a finite set $S \subseteq \X$ if every binary
labeling of $S$ is realized by some concept in $\C$.
\end{definition}

\begin{definition}[VC dimension]
The VC dimension of $\C$, written $\vc(\C)$, is the largest size of a finite set
that $\C$ shatters.
\end{definition}

The growth function counts how many labelings a class can realize on a sample
of size $m$. Sauer's lemma is the bridge between finite VC dimension and
generalization: when $\vc(\C)=d$, the growth function is at most
\[
  \sum_{i=0}^{d} \binom{m}{i}.
\]

The repository code uses this only on small finite examples. The goal is to make
terms like ``shattered'' and ``growth function'' directly checkable, not to
build a symbolic theorem prover for the whole course.
